
\begin{table}[htbp]
   \caption{\label{tab:6yr_huc02fe_inorg_val_sumstats_ravalli_2005} Summary statistics: mean concentration for inorganic chemicals and cumulative upstream coal production}
   \bigskip
   \centering
   \begin{adjustbox}{width = 0.9\textwidth, center}
      \begin{tabular}{lp{3.8cm}ccccc}
         \toprule
         Variable & MCL & Mean & Max & Std.\ Dev. & $N$ & Near MCL \\
         \midrule
         Arsenic (mg/L) & Pre-2006: 0.050 mg/L, 2006+: 0.010 mg/L & 0.0029 & 0.0110 & 0.0023 & 378 & 0.0\% \\
         Nitrate (mg/L) & 10.0 mg/L & 0.7520 & 5.540 & 0.8482 & 444 & 0.5\% \\
         Barium (mg/L) & 2.000 mg/L & 0.0748 & 0.7071 & 0.0671 & 369 & 0.0\% \\
         Chromium (mg/L) & 0.100 mg/L & 0.0059 & 0.0354 & 0.0050 & 374 & 0.0\% \\
         Selenium (mg/L) & 0.050 mg/L & 0.0048 & 0.0354 & 0.0051 & 367 & 0.5\% \\
         Cumul. upstream coal prod. (10M ST) & --- & 0.4395 & 14.179 & 1.415 & 519 & --- \\
         \bottomrule
      \end{tabular}
   \end{adjustbox}
   {\tiny\linespread{1}\selectfont \par \raggedright \textit{Notes:} Sample period 1998--2005. For cumulative upstream coal production, statistics are computed over unique CWS$\times$year pairs that appear in at least one regression sample. ``Near MCL'' is the share of CWS-year mean concentrations exceeding 50\% of the applicable MCL. Sample: community water systems strictly downstream of a coal mine.}
\end{table}

